\documentclass{article}
\usepackage{graphicx,amsmath,amsfonts,amsthm} % Required for inserting images
\newtheorem*{claim}{Claim}

\title{Solutions to Kesavan}
\author{Aditya Padiyar}
\date{February 2026}

\begin{document}

\maketitle

\section{Introduction}
This document contains solutions to \textit{Functional Analysis} by Kesavan

\section*{Chapter 2}

\begin{enumerate}
    \item  It will suffice to prove $c$ is a closed subspace of $\ell^\infty$. The other assertion follows since limit is bounded linear.

    Let $x \in \overline{c}$. Let $\epsilon>0$ and get $y \in C$ with $||x-y||<\frac{\epsilon}{4}$. Now $|y_m-y_n|<\frac{\epsilon}{2}$ for $m,n \geq N$. So $|x_m-x_n|\leq|x_m-y_m|+|x_n-y_n|+|y_m-y_n|<\epsilon$ for $m,n \geq n$

    \item Let $x \in c_0$. Define $x^i=(x_1,x_2,..,x_i,0,0,..) \in c_{00}$. Now $||x-x^i||\to 0$
    \item Follows from Problem 2 since $\ell^1 \supseteq c_{00}$
    \item This is Minkowski's inequality (see MIRA)
    \item 
    \item We know that $C_b(\mathbb{R})$, the space of continuous bounded functions, is a real Banach space. The fact that $C_c(\mathbb{R})$ is a subspace follows from $\text{support }f+g=\overline{(f+g)^{-1}(\mathbb{R}\setminus\{0\} )} \subseteq \overline{f^{-1}(\mathbb{R}\setminus \{0\}) \cup g^{-1}(\mathbb{R}\setminus \{0\})}=\text{support }f \cup \text{support }g$

    Define $f_n(x)=\begin{cases}
        e^{-x^2} \text{ if } |x| \leq n \\
        e^{-n^2}(n+1-|x|) \text{ if } |x| \in (n,n+1) \\
        0 \text{ else}
    \end{cases}$

    This provides a sequence in $C_c(\mathbb{R})$ which converges to $x \to e^{-x^2}$ which is not in $C_c(\mathbb{R})$. Hence $C_c(\mathbb{R})$ is not closed and thus not complete

    \item It is clear that $C_0(\mathbb{R})$ is a subspace of $C_b(\mathbb{R})$. We need to show that is closed.

    Suppose $f_n \to f$ with each $f_n \in C_0(\mathbb{R})$. Let $\epsilon > 0$, there is $N$ such that $||f_n-f|| \leq \frac{\epsilon}{2}$ for $n \geq N$. Get a compact set $K$ such that $|f_N(x)|<\frac{\epsilon}{2}$ outside $K$. Now $|f(x)| \leq |f_N(x)|+|f(x)-f_N(x)| < \epsilon$ outside $K$

    Let $f \in C_0(\mathbb{R})$ and $\epsilon > 0$. We have a compact set $K$ such that $|f(x)|< \frac{\epsilon}{3}$ outside $K$. Put $K \subseteq [-N,N]$. 
    
    Define $g(x)=\begin{cases}
        f(x) \text{ if } |x| \leq N+1 \\
        f(N+1)(N+2-x) \text{ if } x \in (N+1,N+2) \\
        f(-N-1)(N+2+x) \text{ if } x \in (-N-2,-N-1) \\
        0 \text{ else }
    \end{cases}$

    Then $g \in C_c(\mathbb{R})$ and $|f(x)-g(x)| \leq |f(x)|+|g(x)|<\frac{2\epsilon}{3}$ which shows that $C_c(\mathbb{R})$ is dense.

    \item 

    \item Define a norm on $\mathcal{C}^k[0,1]$ by $||f|| = \text{max}_{t \in [0,1]}\{|f^{(0)}(t)|,|f^{(1)}(t)|,..|f^{(k)}(t)|\}$ and suppose $f_1,f_2,..$ is a Cauchy sequence. Note that since $||h|| \geq ||h^{(r)}||_\infty$ we get that $f_1^{(r)},f_2^{(r)},..$ is a Cauchy sequence in the Banach space $C[0,1]$. 
    
    Hence by Theorem 7.17 of Rudin this converges to $f^{(r)}$, for some $f \in \mathcal{C}^k[0,1]$. 
    
    Then by the equation $||h|| \leq \sum_{r=0}^k||h^{(r)}||_\infty$ we get that $f_i \to f$

    \item 
    \item 
    \item 

    \item It suffices to show $\pi(B(0,1))$ is open in $V/W$. Suppose $||v||<1$. 

    \begin{claim}
        $v+W$ is an interior point of $\pi(B(0,1))$
    \end{claim}
    
    \begin{proof}
    Suppose $||x-v+W||<1-||v||$. Then we get $||x-v+w||<1-||v||$. Now $||x+w|| \leq ||x-v+w||+||v||<1$. We also have $\pi(x+w)=x+W$
    \end{proof}

    \item \begin{itemize}
        \item Suppose $T$ is continuous and $v_1,v_2,..$ is a Cauchy sequence. Now $||Tv_n-Tv_m|| \leq ||T|| \ ||v_n-v_m||$ which shows the forward direction.

        \item     Suppose Cauchy sequences are mapped to Cauchy sequences. 
        
        Let $x_n \to 0$. Then $x_1,0,x_2,0,.. \to 0$. We get that $Tx_1,0,Tx_2,0,..$ is Cauchy and hence $||Tx_1||,0,||Tx_2||,0,..$ is Cauchy since norm is uniformly continuous. Hence $Tx_n \to 0$ 
    \end{itemize}
    
    
    


    \item \begin{itemize}
        \item  It is clear that $||f||_\infty+||f'||_\infty \geq ||f||_1 \geq ||f'||_\infty$ which shows that $||T|| \leq 1$

        \item   Now consider $f_n(x)=x^n$. We have $||f_n'||_\infty=n \leq ||T||||f_n||\leq (n+1)||T||$. This shows that $||T|| \geq 1$
    \end{itemize}
    
   
    
  

    \item \begin{itemize}
        \item  We have $T^n1=(x \to \frac{x^n}{n!})$, so $||T^n|| \geq \frac{1}{n!}$

        \item   Now, apply Theorem 5.15 of Rudin (Taylor's theorem) to $T^nf$ to get $T^nf(x)=\frac{f(s)}{n!}x^n$ from which we get $||T^n|| \leq \frac{1}{n!}$
    \end{itemize}
    
    \item 
    \item 
    \item
    \item
    \item
    \item
    \item
    \item
    \item
    \item
    \item
    \item It is easy to see that $T \in \overline{\text{GL}_n(\mathbb{C})}$ for any upper triangular matrix. Since conjugation is a homeomorphism, density follows. Openness follows from 2.39.

    \begin{claim}
        GL$_n(\mathbb{C})$ is path connected
    \end{claim}
    \item
    \item
    \item
    \item
    

    \item We have $||D_w(x)||^2=\sum|w_ix_i|^2 \leq ||w||||x||^2$ from which we get continuity.

    For this map to be invertible the map $x \to (\frac{x_1}{w_1},\frac{x_2}{w_2},..)$ must be a bounded linear operator on $\ell_2$.

    Suppose this is true, and let $C$ be the operator norm. Since $e_i \to \frac{1}{w_i}e_i$ we get $|w_i| \geq C$

    Suppose $w_i$ don't get arbitrarily close to $0$. Then \[||(\frac{x_1}{w_1},\frac{x_2}{w_2},..)||=\sum |\frac{x_i}{w_i}|^2 \leq C^2 ||x||^2\] as required.

    \item
    \item

    \item By Proposition 2.3.3 $\mathcal{L}(V)$ is a Banach space and then use 6.41 of MIRA

    \item We have $\sum_{k=1}^\infty ||A^k|| \leq \sum_{k=1}^\infty ||A||^k<\infty$ where the first inequality follows from 2.3.4. Use Problem 36 to establish convergence. Let $S$ be the sum. Now $||S(I-A)-I||=\text{lim}||S_n(I-A)-I||=\text{lim}||A^{n+1}||=0$. 
    
    A similar calculation on the other side proves the claim.

    \item \begin{enumerate}
        \item Follows from Problem 36
        
        \item Define $E_n(T)=I+\sum_{k=1}^n\frac{T^k}{k!}$. Then $E_n(A+B)-E_n(A)E_n(B)=0$ by the binomial theorem and commutativity, then take limits
        
        \item yes, its inverse is $\text{exp}(-A)$, by (b) and the fact that $A$ commutes with $-A$
        
        \item $\text{exp}(tA)=\text{exp}(t\alpha I+t\omega R_{\frac{\pi}{2}})=e^{\alpha t}\text{exp}(t\omega R_{\frac{\pi}{2}})$, by (b) 
        
        Now use Corollary 3 of Section 1.3 of Perko.
    \end{enumerate} 

    \item From Problem 37 it follows that $I \in \text{interior }\text{GL}(V)$. Let $A$ be invertible. Since the map $B \to AB$ is a homeomorphism and $GL(V)$ is invariant, the result follows.

    \item 
    \begin{enumerate}
        
        \item It is easy to see that $||T||=||S||=1$
        \item 
        \begin{itemize}
            \item We have $||SA-I||=||SA-ST|| <1$. From Problem 37 we have that $SA$ is invertible. Thus $A$ cannot be invertible since $S$ is not.

            \item  Since this ball around $T$ does not intersect $\mathcal{G}$, $\mathcal{G}$ is not dense.
        \end{itemize}
        
    \end{enumerate}

    \item 
    \item 
    \item 
    
    \item \begin{enumerate}
        \item Since $\frac{f(x)}{f(x_0)}x_0 - x\in Z$, we have $||x+Z||=||\frac{f(x)}{f(x_0)}x_0+Z||$ which simplifies to the RHS
        \item see (c)
        \item By (a) we need to prove sup$_{||x|| \leq 1}||x+Z||=1$
        
        Since the map $\pi$ has norm 1 it suffices to show that $||x+Z||>(1-\epsilon)||x||$ for some $x$, and this is shown by 2.3.1.
    \end{enumerate}
\end{enumerate}

\section*{Chapter 3}

\begin{enumerate}
    \item 
    \item 
    \item 
    
    \item Take $x \to x$ and $x \to x^2$ to see that the definition of strict convexity is not satisfied.
    
    \item \begin{itemize}
        \item     Suppose $f \in W^*$ with $f_1 \neq f_2 \in V^*$ extensions such that $||f_1||=||f_2||=||f||$. Now $||\frac{f_1}{||f||}||=||\frac{f_2}{||f||}||=1$ so by strict convexity $||f_1+f_2||<2||f||$. Now $||2fw||=||f_1w+f_2w|| \leq ||f_1+f_2|| \ ||w||$ which is a contradiction

        \item     Suppose $||f_1||=||f_2||=1$ with $f_1(x_0)=f_2(x_0)=||x_0||$. Then we get $f_1=f_2$ by applying the previous result with $f(cx_0)=c||x_0||$
    \end{itemize}

    \item By Corollary 3.1.1 get $\tilde{f} \in (V/\overline{W})^*$ with $||\tilde{f}||=1$. Then $\tilde{f} \circ \pi$ is the required functional.

    \item Let $x_1,..,x_n$ be a basis of $X$ and let $f_1,..,f_n$ be the dual basis. Let $g_i=f_i \circ T$ and by Hahn-Banach get extensions $\tilde{g_i} \in V^*$. Then $v \to \sum (\tilde{g_i}v)x_i$ is the desired extension.

    \item By the previous problem get a map $T \in \mathcal{L}(V,W)$ which is identity on $W$. 
    
    
    ker $T$ is the required subspace since $v-Tv \in \text{ker }T$ 

    ker $T$ is closed since $T$ is continuous.

    If $w \in \text{ker }T$, then $Tw=w=0$. Hence this sum is direct.

    \item \begin{enumerate}
        
        \item for $x \in \ell_1$ define $f_x \in (c_0)^*$ by $f_x(y)=\sum y_i\overline{x_i}$. It is clear that this map is linear, injective and norm preserving.

        let $f \in (c_0)^*$. 
        \begin{claim}
            $x=(\overline{fe_i}) \in \ell_1$
        \end{claim}

        \begin{proof}
            Define $a^n_i=\begin{cases}
                \overline{\frac{fe_i}{|fe_i|}} \text{ if }fe_i \neq 0 \text{ and } i \leq n \\
                0 \text{ else}
            \end{cases}$
            
            Then $fa^n=\sum_{i=1}^n|fe_i|$ 
        \end{proof}

        It is clear that $f$ and $f_x$ agree on $e_i$ and we are done by Problem 2.2

        \item see MIRA 7.26
        \item In $\ell_1$, let $W=\overline{\{\text{span }e_k\}_{k > N}}$
        
        From (a), we have $\ell_1/W \cong (c_0)^*/\overline{\{\text{span }fe_k\}_{k > N}}$

        We have the map $\text{drop}:\ell_1 \to \ell_1^N$ that takes the first $N$ slots. It has kernel $W$. Now for $x \in \ell_1$ we have $||x+W||=||\text{drop }x||$

        In $c_0$, let $H=\text{span}(e_1,..,e_N)$. By Problem 13, we have $H^* \cong (c_0)^*/H^\perp$.

    \end{enumerate}
    
    
    \item 
    \item Suppose $x \in \ell_1$. Get $n$ such that $|x_n| \geq |x_k|$ for all $k$. Now $\sum_{k=1}^N|x_k|(1-\frac{1}{k}) \leq ||x||_1 -\frac{|x_n|}{n}$ for all sufficiently large $N$. 
    
    This shows that $|f_y(x)| \leq ||x||_1-\frac{|x_n|}{n}$. Since $||f_y||=1$ we are done

    \item $p(x)=\text{inf}\{\alpha > 0: ||\frac{x}{\alpha}||<r\}=\frac{||x||}{r}$

    \item \begin{enumerate}
        \item $W^\perp= \cap \text{ker}J_x$ which is an intersection of closed sets
        
        \item \begin{itemize}
            \item Let $g \in W^\perp$ and $||w||\leq 1$. Then $||f-g|| \geq ||(f-g)w||=||fw||$. Take the supremum over $w$ and then the infimum over $g$. This proves LHS$\geq$RHS

            \item Extend $f_{|W}$ to $\tilde{f}$ by Hahn-Banach such that $||\tilde{f}||=||f_{|W}||$. Then LHS$\leq||f-(f-\tilde{f})||=$ RHS where the inequality follows since $f-\tilde{f} \in W^\perp$
        \end{itemize}
        
        \item \begin{itemize}
            \item Suppose $f \in W^*$ and $\tilde{f_1},\tilde{f_2}$ are extensions, then clearly $\tilde{f_1}-\tilde{f_2} \in W^\perp$ which shows that this map is well defined.

            \item Also $||\tilde{f}+W^\perp||=d(\tilde{f},W^\perp)=||f||$ by (b). 

            \item Linearity, injectivity and surjectivity are trivial to check.
        \end{itemize}
        
        \item \begin{itemize}
            \item Suppose $f \in W^\perp$. It is easy to check that $v+W \to fv$ is a well defined map. We claim that this map is in the pre image of $f$.
            
            Now $||fv||=||f(v+w)|| \leq ||f|| \ ||v+w||$. Hence $||fv|| \leq ||f|| \ ||v+W||$ which shows that this map has norm at most $||f||$
            
            This shows surjectivity.

            \item Injectivity is easy to check. 
            
            \item \begin{itemize}
                \item $||\tau g|| = ||g \circ\pi|| \leq ||g||$ since $||\pi||=1$. 

                \item The other inequality is shown by the first bullet point, since $\tau$ is injective.
            \end{itemize}
            
            
            
            
        \end{itemize}
    \end{enumerate}

    \item \begin{enumerate}
        \item $Z^\perp=\cap \text{ker}g$
        \item Suppose $x_0 \in (W^\perp)^\perp$ with $x_0 \notin \overline{W}$. By Hahn-Banach get $f \in W^\perp$ such that $f(x_0) \neq 0$ 
        \item trivial
        \item Suppose $f \in (Z^\perp)^\perp$ but $f \notin \overline{Z}$. By Hahn-Banach get $J_x \in V^{**}$ such that $J_xZ=0$ but $J_xf \neq 0$  
    \end{enumerate}
    
    \item follows directly from definition
    \item follows directly from definition, since $fA \in V^*$ for any $f \in W^*$
    \item follows from definition and the corresponding property for scalar valued functions
    \item trivial
    \item \begin{enumerate}
        \item $f \to f(s)$ is a bounded linear functional
        \item follows from Fubini's theorem
    \end{enumerate}
    \item We need to show $\frac{\int_t^{t+h}[\phi(s)-\phi(t)]ds}{h} \to 0$. Let $\epsilon > 0$ and get $\delta>0$ such that $||\phi(s)-\phi(t)||<\epsilon$ if $|s-t|<\delta$. Then the norm of the fraction is bounded by $\frac{\int_t^{t+h}||\phi(s)-\phi(t)||ds}{h}$. If $h<\delta/2$ we can bound this by $\epsilon$
    \item 
    \item 
    
    
    \item Define $A:V\to \mathbb{F}^n$ by $A=f_1 \times .. \times f_n$. 
    
    Now ker$A \subseteq \text{ker }f_0$, so $f_0=\tilde{f_0} \circ \pi$. 
    
    Similarly, ker$A \subseteq \text{ker}f_i$ so $f_i=\tilde{f_i}\circ \pi$. 
    
    By Problem 3F.25 of LADR, $\tilde{f_0}=\sum\alpha_i\tilde{f_i}$. Composing with $\pi$ finishes the job.
    
\end{enumerate}

\section*{Chapter 4}

\begin{enumerate}
    \item Suppose $v_1,v_2,..$ is a basis for a Banach space $V$. 
    
    Then $V=\bigcup_{i}\text{span}(v_i,..,v_i)$. This violates Baire's theorem.

    \item \begin{enumerate}
        \item This follows from 2.28
        \item it is easy to see that $0 \in \overline{\text{NonZeroTrace}}$. Now add a zero trace matrix to both sides.
    \end{enumerate}

    \item \begin{enumerate}
        \item We have \[(x \to \sum_{i=1}^na_ix_i) \in (c_0)^*\]
              By Corollary 4.2.2 $(x \to \sum a_ix_i) \in (c_0)^*$. Now use Problem 3.9

        \item We have $(x \to \sum_{i=1}^na_ix_i) \in (\ell_p)^*$ by Holder's inequality. Now the same technique works.
    \end{enumerate}

    \item 
    \item 
    \item 
    \item Suppose $x_n \to x$ and $Tx_n \to y$. We have $a(Tx_n,Tx-y)=a(x_n,T^2x-Ty)$, from which we see that $a(y,Tx-y)=a(x,T^2x-Ty)=a(Tx,Tx-y)$ and now use the closed graph theorem.

    \item Suppose $x_n \to x$ and $Tx_n \to y$. Suppose $Tx-y \neq 0$ so get $f_j$ such that $f_jTx-f_jy \neq 0$. But we have $f_jTx_n \to f_jTx$ and $f_jTx_n \to f_jy$

    \item Suppose $x_n \to x$ and $Bx_n \to y$. Then $ABx_n=Tx_n \to Ay=Tx$. So $Bx=y$

    \item Suppose $f \in W^\perp$. Let $w_n \to w \in \overline{W}$. Now $fw_n \to fw=0$

    \item Suppose $g \in (\text{ker}f)^\perp$. Assume $f(x_0)=1$. Now $x-f(x)x_0 \in \text{ker}f$
    \item the map lim has a right inverse which is the constant sequence

    \item Suppose $ST=I$. Then $w-TSw \in \text{ker}S$, which shows ker $S \oplus$range $T=W$. Suppose $Tx_n \to w$. Then $x_n \to Sw$ and $Tx_n \to TSw$

    Suppose $W=T(V) \oplus N$ and let $P$ be the first projection. Define $S$ by $TSw=Pw$. Now $T(STv-v)=PTv-Tv=0$. Suppose $w_n \to w$ and $Sw_n \to v$. Then $TSw_n=Pw_n \to Pw=Tv=TSw$. Hence $T(v-Sw)=0$

    \item \begin{enumerate}
        \item 
            Let $x_n$ be Cauchy in $V$. Then $Tx_n$ and $x_n$ are Cauchy in $W$. So $x_n \to w_1$ and $Tx_n \to w_2$. Since the graph is closed, $w_1 \in V$ and $w_2=Tw_1$. Now $||x_n-w_1||=||x_n-w_1||_W+||Tx_n-Tw_1||_W \to 0$
        \item 
        Get constants $C_1$ and $C_2$ such that $||x||_W \leq C_1||x||_V$ and $||Tx||_W \leq C_2||x||V$. Adding, we get that the identity map on $V$ is a homeomorphism using the open mapping theorem.
    \end{enumerate}

    \item use the open mapping theorem

    \item For $z \in \text{ker}T$, we have $||T(v+z)|| =||Tv|| \geq c||v+z||$ which shows continuity of $(\tilde{T})^{-1}$,and the open mapping theorem finishes the job.

    \item $T$ is invertible, $\iff T$ is bijective and continuous, $\iff T^*$ is bijective and continuous, $\iff T^*$ is invertible

    To get the second equivalence use Theorem 4.7.1 , Corollary 4.7.1, and Proposition 4.7.3

    \item follows from Corollary 4.7.1
\end{enumerate}

\section*{Chapter 5}
\begin{enumerate}
    \item Suppose a topology is such that all $f \in W^*$ are continuous. For $U \subseteq \mathbb{F}$ open and $f \in V^*$, we have $i^{-1}f^{-1}U=(f \circ i)^{-1}U$ which is open.

    Suppose a topology is such that $i:W \to V_{\text{weak}}$ is continuous. Let $f \in W^*$. Extend $f$ to $g \in V^*$. Let $U \subseteq \mathbb{F}$ be open. Now $f^{-1}U=(g^{-1}U) \cap W=i^{-1}g^{-1}U$ which is open.

    \item \begin{itemize}
        \item Since $\overline{W}$ is convex and norm closed, from Proposition 5.1.4 it is weakly closed. This shows that that it contains the weak closure.

        \item The other inclusion follows since the norm topology is finer than the weak topology.
    \end{itemize}
    
    \item Suppose $f$ is in the weak* closure of $W^\perp$. 
    
    Let $w\in W$. A weak* neighborhood of $f$ is $\{g:|(f-g)(w)|<\frac{1}{n}\}$. 
    
    Hence get $g_n \in W^\perp$ with $g_n(w) \to f(w)=0$

    \item Consider the weak open set $U=\{x:|f_i(x-x_0)|<\epsilon\}$. Let $x,y \in U$. Then $|f_i(tx+(1-t)y-(t+1-t)x_0)| < \epsilon$ by the triangle inequality

    Consider the weak* open set $U=\{f:|(f-f_0)(x_i)|<\epsilon\}$. Let $f,g \in U$. Then $|(tf+(1-t)g-(t+1-t)f_0)(x_i)|<\epsilon$

    \item 
    \item 
    \item 
    \item 
    \item \begin{claim}
        in $c_0$ we have $e_n\rightharpoonup0$
    \end{claim}

    \begin{proof}
        Fix $x \in \ell_1$. Then $f_xe_n=\overline{x_n} \to 0$. But $(c_0)^*=\ell_1$
    \end{proof}

    \begin{claim}
        in $\ell_\infty$ we have $e_n\rightharpoonup0$
    \end{claim}

    \begin{proof}
        Let $f \in (\ell_\infty)^*$. Since $f_{|c_0} \in (c_0)^*$ we have $fe_n \to 0$ by the previous claim.
    \end{proof}

    \item 
    \item 
    
    \item From Corollary 5.1.1 we have $||x|| \leq \text{liminf }||x_n||$. From the problem we get $||x|| \geq \text{limsup }||x_n||$

    \item 
    \item 
    \item
    \item
    
    \item Assume $V$ is infinite dimensional. Consider a weak* open neighborhood $U=\{f:|f(x_i)|<\epsilon\}$. Apply Corollary 3.3.1 to $W=$ span $x_i$ to get $0 \neq f \in W^\perp$. We can scale this functional to an arbitrarily large norm and it is in $U$
    
    \item 
    \item 
    \item 
    \item 
    
    \item \begin{itemize}
        \item We have $||Tx|| \leq \frac{\pi^2}{6}||x||$ which shows continuity
        
        \item $T(e_1+..+e_n)=(1,\frac{1}{4},..,\frac{1}{n^2},0,0,..) \to \{\frac{1}{n^2}\} \notin \text{range }T$. This shows that $T(B)$ is not closed.
    \end{itemize}
    
    \item From Theorem 5.3.1 $B$ is weakly compact. Hence, $T(B)$ is weakly compact and thus closed.
    
    \item Suppose $V \subseteq \ell_1$ is reflexive.
    \begin{claim}
        $B_V$, the closed unit ball in $V$, is compact
    \end{claim}

    \begin{proof}
        Suppose $v_n \in B_V$. By Theorem 5.4.2 $v_{n_k} \rightharpoonup v$. By Schur's lemma $v_{n_k} \to v$. Since $B_V$ is closed $v \in B_V$. We have shown sequential compactness.
    \end{proof}

    By Proposition 2.3.6 $V$ is finite dimensional
    
    \item We have $c_{00} \subseteq \ell_p \subseteq c_0$. It is easy so see that $S=\bigcup_{n \geq 1}\mathbb{Q}^n$ is dense in $c_{00}$. By Problem 2.2 $S$ is dense in $c_0$. Since on $\ell_p$ the sup norm is finer than the p-norm, we are done.
    
    \item See the previous problem
    
    \item By Stone-Weierstrass, polynomials are dense in $\mathcal{C}[0,1]$. Separability follows since any polynomial can be approximated with one with rational coefficients, or, if the field is $\mathbb{C}$, one with coefficients having rational real and imaginary parts.

\end{enumerate}


\section*{Chapter 6}

\begin{enumerate}
    \item 
    \item
    \item
    \item
    \item
    \item
    \item
    \item
    \item
    \item
    \item
    \item
    \item
    \item
    \item
    \item
    \item
    \item
    \item
    \item
    \item
    \item
    \item
    \item
    \item
    \item
    \item
\end{enumerate}

\section*{Chapter 7}

\begin{enumerate}
    \item 
    \item 
    \item 
    \item 
    \item 
    \item 
    \item 
    \item 
    \item 
    \item 
    \item 
    \item 
    \item 
    \item 
    \item 
    \item 
    \item 
    \item 
    \item 
    \item 
    \item 
    \item 
    \item 
    \item 
    \item 
    \item 
    \item 
    \item 
    \item 
    \item 
    \item 
    \item 
    \item 
    \item 
    \item 
    \item 
    \item 
\end{enumerate}

\section*{Chapter 8}

\begin{enumerate}
    \item By Proposition 5.1.2 $x_n$ is bounded. Suppose $T(x_n)$ does not converge to $T(x)$. Then we have $n_k$ such that $||T(x_{n_k})-T(x)|| \geq \epsilon$. By Remark 8.1.1 $T(x_{n_{k_j}}) \to w$. Since $T$ is also weakly continuous we have that $T(x)=w$ which is a contradiction.
    
    \item \begin{enumerate}
        \item Suppose $x_n$ is bounded. By Theorem 5.4.2 $x_{n_k} \rightharpoonup x$. By hypothesis $T(x_{n_k}) \to T(x)$. By Remark 8.1.1 $T$ is compact.

        \item Consider the identity operator on $\ell_1$. By Schur's lemma it satisfies the given property. But by Example 8.1.3 it is not compact.
    \end{enumerate}
    
    \item In the Hilbert space $\ell_2$ let $M=\overline{\text{span}(e_1,e_3,..)}$. We have that $P_MP_{M^\perp}=0$ which is compact. But by 10.74 of MIRA neither $P_M$ nor $P_{M^\perp}$ is compact.
    
    \item 
    
    \begin{enumerate}
        
        \item This follows from Problem 2 (a) and the Schur property of $\ell_1$

        \item We have $T^* \in \mathcal{L}(W^*,(c_0)^*)$. By Problem 3.9, we have an isometric isomorphism $\phi:(c_0)^* \to \ell_1$. By Corollary 5.3.3 and Problem 2 (a), $\phi \circ T^*$ is compact. Hence, $T^*$ is compact, which proves the claim by Proposition 8.1.3.
    
    \end{enumerate}
    
    \item Assume that $T$ is injective. Then, by Corollary 4.4.1, we have $T^{-1} \in \mathcal{L}(\mathcal{R}(T),c_0)$. By the previous problem, $T$ is compact. It remains compact when we restrict the co-domain to $\mathcal{R}(T)$. This contradicts Example 8.1.4.
    
    \item This follows from 10.74 of MIRA
    
    \item Linearity is obvious. Using the Cauchy-Schwarz inequality, it is easy to show that the map is bounded. Since it has finite rank, it is compact.
    
    \item By the Riesz Representation Theorem, get $v_k \in H$ such that $\langle x,v_k \rangle = \langle Tx, w_k \rangle$ for all $x \in H$
    
    \item We have that $T \in B(H)$ by Problem 7. It is easy to see that $T^2=0$. Therefore, we have $(I+\alpha T)(I - \alpha T)=I$ which proves the claim.
    
    \item Assume $\langle Ax,x \rangle \geq \alpha ||x||^2$ for all $x \in H$, with $\alpha>0$. Find an orthonormal sequence $e_1,e_2,..$. By Bessel's inequality (Theorem 7.4.1), we have that $e_n \rightharpoonup 0$. By Problem 1, we have that $Te_n \to 0$. By the hypothesis and the Cauchy-Schwarz inequality, we have $\alpha \leq \langle Te_n,e_n \rangle  \leq ||Te_n||$. Taking the limit, we get a contradiction.
    
    \item 
    
    \begin{itemize}
        \item Suppose $S(I-T)=I$. Then $I-T$ is injective. By Theorem 8.2.1, $I-T$ is surjective as well. Then $S=(I-T)^{-1}$, which proves the forward direction.

        \item Suppose $(I-T)S=I$. Then $I-T$ is surjective. By Theorem 8.2.1, $I-T$ is injective as well. Then $S=(I-T)^{-1}$, which proves the backward direction.
    \end{itemize}

    Suppose $S=(I-T)^{-1}$. Then we have $I-S=-ST$. By Example 8.1.4, $-ST$ is compact, and hence so is $I-S$
    
    \item We have $||Tx_n|| \to ||T||$ where $||x_n||=1$. By Theorem 5.4.2, we have $x_{n_k} \rightharpoonup x$. By Problem 1 and the continuity of the norm $||Tx_{n_k}|| \to ||Tx|| = ||T|| \leq ||T|| ||x||$. Dividing, we get $||x|| \geq 1$ and by Proposition 5.1.2 we get $||x|| \leq 1$.
    
    \item Suppose, for the sake of contradiction, that $||u_n||_2> \epsilon ||u_n||_1+n||u_n||_3$. We may assume that $||u_n||_1=1$. By Theorem 5.4.2 we get $u_{n_k} \rightharpoonup u$ in $H_1$. By Problem 1 and the continuity of the inclusions, $u_{n_k} \to u$ in $H_2$ and $H_3$. We have $||u_{n_k}||_2 > \epsilon + n_k ||u_{n_k}||_3$ from which it follows that $u=0$. Then it follows that $\epsilon \leq 0$, which is a contradiction.
    
    \item Define the linear map $\text{Hankel}:\ell_1 \to \mathcal{L}(\ell_2)$ by \[(\text{Hankel}_ax)_k=\sum_{j=1}^\infty x_ja_{j+k-1}\]
    
    \begin{claim}
        $||\text{Hankel}|| \leq 1$ 
    \end{claim}

    \begin{proof}
        Let $a \in \ell_1$ and $b=(|a_1|,|a_2|,..) \in \ell_1$. It is easy to see that $||\text{Hankel}_a|| \leq ||\text{Hankel}_b||$. By Schur's test on page 41 of the textbook, we have $||\text{Hankel}_b|| \leq ||a||_1$
    \end{proof}

    \begin{claim}
        $\mathcal{R}(\text{Hankel}) \subseteq \mathcal{K}(\ell_2)$
    \end{claim}

    \begin{proof}
        We know that $c_{00}$ is dense in $\ell_1$. It is easy to see that $\text{Hankel}(c_{00}) \subseteq \text{FiniteRank}(\ell_2) \subseteq \mathcal{K}(\ell_2)$. By Theorem 18.1 of Munkres we have \\ $\mathcal{R}(\text{Hankel}) \subseteq \overline{\text{Hankel}(c_{00})}$. Proposition 8.1.1 completes the proof.
    \end{proof}
    
    \item 
    \item 
    \item 
    \item 
    \item 
    \item 
    \item 
    \item 
    \item 
    \item 
    \item 
    \item 
    \item 
    \item 
    \item 
    \item 
    \item 
    \item 
    \item 
\end{enumerate}

\end{document}
